% GENERATED FILE. Edit canonical lessons, then run npm run generate:books.
% canonical-source-hash: fnv1a64:9cbee480ed5f1a6e
% canonical-lessons: JA-C399-makura, JA-C399-moufu, JA-C399-zabuton, JA-C399-mezameru, JA-C399-nekorobu

\chapter{Futon, Pillow and Blanket}
\label{ch:ja-futon-pillow-and-blanket}
\hlchaptermodality{399}

\begin{chapteropening}
\textbf{By the end of this chapter:} \emph{I can ask for bedding and get ready for bed.}

Complete the last lesson of chapter 399: I can ask for bedding and get ready for bed.
\end{chapteropening}

\section[\texorpdfstring{\ja{まくら} (makura)}{まくら (makura)}]{\ja{まくら} (makura) — a pillow}

\label{lesson:JA-C399-makura}



\begin{quote}
Before the new one: say the Japanese for to soak, then the Japanese for to relax.
\end{quote}

\subsection*{You'll want to know: \ja{まくら}}
\textbf{\ja{まくら}} — \emph{makura} — \textquotedblleft{}a pillow\textquotedblright{}.

\begin{quote}
\textbf{\ja{この} \ja{まくらは} \ja{かたいです。}} — \emph{kono makura wa katai desu} — \textquotedblleft{}This pillow is hard.\textquotedblright{}
\end{quote}

\begin{cousinweb}
\textbf{\ja{まくら}} (\emph{makura}) is a native Japanese word (wago), inherited rather than borrowed from Chinese; its kanji was attached later as a kun'yomi (JMdict).
\end{cousinweb}

\begin{grammarlens}[title={What you've built}]
One more word for getting around.
\end{grammarlens}

\subsection*{Your turn}
\begin{itemize}
\raggedright
  \item \emph{Say it:} \emph{makura}
  \item \emph{Say it:} \emph{makura}, once more
  \item \emph{Say it:} \emph{kono makura wa katai desu} — This pillow is hard.
  \item \emph{Recall:} say the Japanese for to soak, then the Japanese for to relax, then say \emph{makura} again
\end{itemize}

\begin{tcolorbox}[breakable,colback=teal!4,colframe=teal!35!black,title={Before you move on}]
What does \ja{まくら} mean? (\textquotedblleft{}A pillow\textquotedblright{}.) Say it once more.
\end{tcolorbox}
\section[\texorpdfstring{\ja{もうふ} (mōfu)}{もうふ (mōfu)}]{\ja{もうふ} (mōfu) — a blanket}

\label{lesson:JA-C399-moufu}



\begin{quote}
Before the new one: say the Japanese for to relax, then the Japanese for a pillow.
\end{quote}

\subsection*{You'll want to know: \ja{もうふ}}
\textbf{\ja{もうふ}} — \emph{mōfu} — \textquotedblleft{}a blanket\textquotedblright{}.

\begin{quote}
\textbf{\ja{さむいので、もうふを} \ja{もう} \ja{いちまい} \ja{ください。}} — \emph{samui node, mōfu o mō ichimai kudasai} — \textquotedblleft{}It is cold, so please give me one more blanket.\textquotedblright{}
\end{quote}

\begin{cousinweb}
\textbf{\ja{もうふ}} (\emph{mōfu}) is Sino-Japanese (kango): two Chinese characters read with their on'yomi, readings borrowed from Chinese. They are \emph{mō}, \textquotedblleft{}fur\textquotedblright{} (Middle Chinese \emph{maw}) and \emph{fu}, \textquotedblleft{}cloth\textquotedblright{} (\emph{puH}) (KANJIDIC2; Middle Chinese in Baxter's transcription).
\end{cousinweb}

\begin{grammarlens}[title={What you've built}]
One more word for getting around.
\end{grammarlens}

\subsection*{Your turn}
\begin{itemize}
\raggedright
  \item \emph{Say it:} \emph{mōfu}
  \item \emph{Say it:} \emph{mōfu}, once more
  \item \emph{Say it:} \emph{samui node, mōfu o mō ichimai kudasai} — It is cold, so please give me one more blanket.
  \item \emph{Recall:} say the Japanese for to relax, then the Japanese for a pillow, then say \emph{mōfu} again
\end{itemize}

\begin{tcolorbox}[breakable,colback=teal!4,colframe=teal!35!black,title={Before you move on}]
What does \ja{もうふ} mean? (\textquotedblleft{}A blanket\textquotedblright{}.) Say it once more.
\end{tcolorbox}
\section[\texorpdfstring{\ja{ざぶとん} (zabuton)}{ざぶとん (zabuton)}]{\ja{ざぶとん} (zabuton) — a floor cushion}

\label{lesson:JA-C399-zabuton}



\begin{quote}
Before the new one: say the Japanese for a pillow, then the Japanese for a blanket.
\end{quote}

\subsection*{You'll want to know: \ja{ざぶとん}}
\textbf{\ja{ざぶとん}} — \emph{zabuton} — \textquotedblleft{}a floor cushion\textquotedblright{}.

\begin{quote}
\textbf{\ja{ざぶとんに} \ja{すわって} \ja{ください。}} — \emph{zabuton ni suwatte kudasai} — \textquotedblleft{}Please sit on the cushion.\textquotedblright{}
\end{quote}

\begin{cousinweb}
\textbf{\ja{ざぶとん}} (\emph{zabuton}) is Sino-Japanese (kango): \emph{za}, \textquotedblleft{}seat\textquotedblright{} (Middle Chinese \emph{dzwaH}), read with its on'yomi, a reading borrowed from Chinese, and the taught \emph{futon}, \textquotedblleft{}bedding\textquotedblright{}, voiced as \emph{-buton}; \emph{futon} is usually traced to a Chinese word for a rush cushion (JMdict, KANJIDIC2; Middle Chinese in Baxter's transcription).
\end{cousinweb}

\begin{grammarlens}[title={What you've built}]
One more word for getting around.
\end{grammarlens}

\subsection*{Your turn}
\begin{itemize}
\raggedright
  \item \emph{Say it:} \emph{zabuton}
  \item \emph{Say it:} \emph{zabuton}, once more
  \item \emph{Say it:} \emph{zabuton ni suwatte kudasai} — Please sit on the cushion.
  \item \emph{Recall:} say the Japanese for a pillow, then the Japanese for a blanket, then say \emph{zabuton} again
\end{itemize}

\begin{tcolorbox}[breakable,colback=teal!4,colframe=teal!35!black,title={Before you move on}]
What does \ja{ざぶとん} mean? (\textquotedblleft{}A floor cushion\textquotedblright{}.) Say it once more.
\end{tcolorbox}
\section[\texorpdfstring{\ja{めざめる} (mezameru)}{めざめる (mezameru)}]{\ja{めざめる} (mezameru) — to wake from sleep; to awaken}

\label{lesson:JA-C399-mezameru}



\begin{quote}
Before the new one: say the Japanese for a blanket, then the Japanese for a floor cushion.
\end{quote}

\subsection*{You'll want to know: \ja{めざめる}}
\textbf{\ja{めざめる}} — \emph{mezameru} — \textquotedblleft{}to wake from sleep; to awaken\textquotedblright{}.

\begin{quote}
\textbf{\ja{りょかんの} \ja{ふとんで} \ja{あさ} \ja{はやく} \ja{めざめました。}} — \emph{ryokan no futon de asa hayaku mezamemashita} — \textquotedblleft{}I woke up early in the morning on the inn's futon.\textquotedblright{}
\end{quote}

Polite \textbf{\ja{めざめます}}, \emph{mezamemasu}; an ichidan verb. Past (polite): \textbf{\ja{こどもが} \ja{よなかに} \ja{めざめました。}}, \emph{kodomo ga yonaka ni mezamemashita}.

\begin{cousinweb}
\textbf{\ja{めざめる}} (\emph{mezameru}) is a native Japanese word (wago), inherited rather than borrowed from Chinese, built from \emph{me}, \textquotedblleft{}eye\textquotedblright{}, and \emph{zameru}, \textquotedblleft{}waking\textquotedblright{} (from \emph{sameru}, \textquotedblleft{}to wake up\textquotedblright{}); its kanji were attached later as kun'yomi (JMdict).
\end{cousinweb}

\begin{grammarlens}[title={What you've built}]
One more word for getting around.
\end{grammarlens}

\subsection*{Your turn}
\begin{itemize}
\raggedright
  \item \emph{Say it:} \emph{mezameru}
  \item \emph{Say it:} \emph{mezameru}, once more
  \item \emph{Say it:} \emph{ryokan no futon de asa hayaku mezamemashita} — I woke up early in the morning on the inn's futon.
  \item \emph{Recall:} say the Japanese for a blanket, then the Japanese for a floor cushion, then say \emph{mezameru} again
\end{itemize}

\begin{tcolorbox}[breakable,colback=teal!4,colframe=teal!35!black,title={Before you move on}]
What does \ja{めざめる} mean? (\textquotedblleft{}To wake from sleep; to awaken\textquotedblright{}.) Say it once more.
\end{tcolorbox}
\section[\texorpdfstring{\ja{ねころぶ} (nekorobu)}{ねころぶ (nekorobu)}]{\ja{ねころぶ} (nekorobu) — to lie down; to stretch out}

\label{lesson:JA-C399-nekorobu}



\begin{quote}
Before the new one: say the Japanese for a floor cushion, then the Japanese for to wake from sleep.
\end{quote}

\subsection*{You'll want to know: \ja{ねころぶ}}
\textbf{\ja{ねころぶ}} — \emph{nekorobu} — \textquotedblleft{}to lie down; to stretch out\textquotedblright{}.

\begin{quote}
\textbf{\ja{たたみの} \ja{うえに} \ja{ねころびました。}} — \emph{tatami no ue ni nekorobimashita} — \textquotedblleft{}I stretched out on the tatami.\textquotedblright{}
\end{quote}

Polite \textbf{\ja{ねころびます}}, \emph{nekorobimasu}; a godan verb. Past (polite): \textbf{\ja{こどもが} \ja{たたみに} \ja{ねころびました。}}, \emph{kodomo ga tatami ni nekorobimashita}.

\begin{cousinweb}
\textbf{\ja{ねころぶ}} (\emph{nekorobu}) is a native Japanese word (wago), inherited rather than borrowed from Chinese, built from \emph{ne}, \textquotedblleft{}sleeping\textquotedblright{} (from \emph{neru}) and \emph{korobu}, \textquotedblleft{}to tumble\textquotedblright{}; its kanji were attached later as kun'yomi (JMdict).
\end{cousinweb}

\begin{grammarlens}[title={What you've built}]
One more word for getting around.
\end{grammarlens}

\subsection*{Your turn}
\begin{itemize}
\raggedright
  \item \emph{Say it:} \emph{nekorobu}
  \item \emph{Say it:} \emph{nekorobu}, once more
  \item \emph{Say it:} \emph{tatami no ue ni nekorobimashita} — I stretched out on the tatami.
  \item \emph{Recall:} say the Japanese for a floor cushion, then the Japanese for to wake from sleep, then say \emph{nekorobu} again
\end{itemize}

\begin{tcolorbox}[breakable,colback=teal!4,colframe=teal!35!black,title={Before you move on}]
What does \ja{ねころぶ} mean? (\textquotedblleft{}To lie down; to stretch out\textquotedblright{}.) Say it once more.
\end{tcolorbox}
